\subsection{Selecting a criterion in normalized smoothness coordinates}
\label{sec:smoothness-coordinate-search}
The monotonicity in Eq.~\eqref{eq:s-derivative} permits a change of
optimization coordinate without changing the underlying PLS estimator
or statistical selection criterion. Write $r=n-d$, let
$\delta_1,\ldots,\delta_r$ denote the strictly positive penalty
eigenvalues, and define
\begin{equation}\label{eq:inverse-s-coordinate}
s=S(\lambda)
=\frac{1}{r}\sum_{j=1}^{r}
\frac{\lambda\delta_j}{1+\lambda\delta_j},
\qquad
\Lambda(s)=S^{-1}(s),\quad 0<s<1.
\end{equation}
Since
\begin{equation}\label{eq:s-coordinate-derivative}
S'(\lambda)=\frac1r\sum_{j=1}^{r}
\frac{\delta_j}{(1+\lambda\delta_j)^2}>0
\qquad (\lambda\ge0),
\end{equation}
$S$ is continuous and strictly increasing, taking the extended
endpoints $S(0)=0$ and $S(\infty)=1$. We define
$\Lambda(0)=0$ and $\Lambda(1)=\infty$ exactly. These formulas
apply for any admissible order $0\le d<n$, provided the positive
penalty spectrum is numerically resolved.

For an existing criterion $C(\lambda;\mu)$ conditional on a fixed
reference mean, introduce $\widetilde C(s;\mu)=C(\Lambda(s);\mu)$.
The same ideal optimization problem can be written as
\begin{equation}\label{eq:criterion-s-coordinate}
\widehat s\in\arg\min_{s\in[0,1]}
\widetilde C(s;\mu),
\qquad
\widehat\lambda=\Lambda(\widehat s).
\end{equation}
Here inadmissible boundary values retain their original definition,
including an undefined ordinary CV score at $\lambda=0$; the
correspondence between global minimizers is set-valued in the
presence of ties. This is an exact reparameterization, not a new
criterion or a proof that a finite numerical search locates the
global optimum. Equal $s$ across difference orders means equal
\emph{normalized conditional trace smoothness}, not equal trend
geometry, variability, or forecast accuracy.

For a differentiable criterion and $0<s<1$, the inverse-function
theorem and the chain rule give
\begin{equation}\label{eq:score-s-chain}
\frac{d\widetilde C}{ds}
=\frac{C'(\lambda)}{S'(\lambda)}
=\frac{dC/d\log\lambda}{dS/d\log\lambda},
\qquad
\frac{dS}{d\log\lambda}
=\frac{1}{r}\sum_{j=1}^{r}
\frac{\lambda\delta_j}{(1+\lambda\delta_j)^2}>0,
\quad\lambda=\Lambda(s).
\end{equation}
Thus interior stationary points coincide in the two
parameterizations. The public implementation exposes the optional
\texttt{search\_space="smoothness"} in
\texttt{select\_lambda}, while
\texttt{search\_space="lambda"} remains the default.
For classical fixed-$\mu$ CV, it scans a uniform and a
logit-spaced set of $s$ values, refines sampled derivative
sign-change brackets and score valleys, and compares evaluated
criterion values at those candidates and exact boundaries.
For GCV, AICc, BIC, and conditional Gaussian MLE the option
evaluates a grid in $s$, using the same score as in $\lambda$.
A user-supplied \texttt{grid} is interpreted in the selected
coordinate. The returned \texttt{lambda\_} remains a penalty;
\texttt{smoothness} returns its normalized trace index,
\texttt{grid} records evaluated penalties, and
\texttt{search\_grid} records the requested coordinates.
Only Gaussian MLE supports the previously defined
observed-difference plug-in for $\mu$; the new coordinate does
not jointly estimate $(\lambda,\mu)$.

The inverse is obtained with a safeguarded root finder in
$\log\lambda$, reusing the spectral representation.
Every interior smoothness evaluation therefore incurs an
additional scalar inversion, typically several $O(r)$
spectral evaluations, before the unchanged criterion score.
Moreover $dS/d\log\lambda$ tends to zero at both limits, so
derivative conversion can be poorly conditioned near $s=0$
or $s=1$. Interior values requiring an unrepresentable
floating-point penalty are rejected rather than silently
treated as endpoints. A bounded search interval does not,
by itself, establish a speed or accuracy advantage over
direct optimization in $\log\lambda$.

For completeness, the prospective joint likelihood profile
in Eq.~\eqref{eq:lambda-profile-mle} could also be expressed
in this coordinate as
\begin{equation}\label{eq:future-joint-s-profile}
\widehat s_{\mathrm{joint}}
=\arg\min_{0<s<1}
J\bigl(\Lambda(s);\widehat\mu(\Lambda(s))\bigr),
\qquad
\widehat\lambda_{\mathrm{joint}}
=\Lambda(\widehat s_{\mathrm{joint}}).
\end{equation}
Equation~\eqref{eq:future-joint-s-profile} describes an
\emph{unimplemented research extension}, not a result of
the conditional-MLE selector or of the benchmarks below.
Boundary candidates and identifiability would require
separate treatment in a full joint-estimation study.
